\slajdid{09_2-s033}
% element: 09_2-s033:e01
% element: 09_2-s033:e02
% element: 09_2-s033:e03
\begin{frame}[label=09_2-s033]{CMOS: izvođenje praga u četvrtoj oblasti}
\begin{columns}[T,onlytextwidth]\column{.24\textwidth}\centering {\RazaviSetup
\begin{circuitikz}[x=1cm,y=1cm,font=\fontsize{10}{12}\selectfont]
\node[rz nmos] (N) at (.28,.65) {};
\node[rz pmos] (P) at (.28,2.2) {};
\coordinate (O) at (N.D |- 0,1.43);
\draw[wire] (N.D)--(O)--(P.D);
\draw[wire] (N.S)--++(0,-.20) coordinate (GND);
\RazaviGround{GND}
\draw[wire] (P.S)--++(0,.22) coordinate (VDD);
\RazaviSupply{VDD}{V_{DD}}
\draw[wire] (O)--++(.60,0) node[rz terminal]{} node[right=3pt]{$V_O$};
\fill (O) circle(1.15pt);
\node[right=3pt] at (N.east) {$\mathrm{N}$};
\node[right=3pt] at (P.east) {$\mathrm{P}$};
\coordinate (I) at (-.7,1.43);
\draw[wire] (N.G)--(I |- N.G)--(I |- P.G)--(P.G);
\draw[wire] (I)--++(-.25,0) node[rz terminal]{} node[left=3pt]{$V_I$};
\fill (I) circle(1.15pt);
\end{circuitikz}}
\column{.72\textwidth}Za $V_{IL}$ koristimo nagib $-1$ i aproksimacije:
\[\frac{V_{O(IH)}-V_{DD}}{L_nE_{Cn}}\ll1,\qquad\frac{V_{IH}-V_{DD}-V_{Tp}}{L_pE_{Cp}}\ll1\]
\[\alert{1.}\quad k_p(V_I-V_{DD}-V_{Tp})^2=k_n\left(2V_O(V_I-V_{Tn})-V_O^2\right)\]
\[\begin{aligned}\alert{2.}\quad V_I&=\frac{2k_nV_O+k_pV_{DD}+k_pV_{Tp}+k_nV_{Tn}}{k_p+k_n}\\&={\color{DEplava}\frac{2\frac{k_n}{k_p}V_O+V_{DD}+V_{Tp}+\frac{k_p}{k_n}V_{Tn}}{1+\frac{k_n}{k_p}}}\end{aligned}\]
\end{columns}
\end{frame}
